\scanpage{219}
% archive sheet number: 218

$X$ \deleted{\unsure{schéma}} préschéma de type fini \inserted{normal} irréductible sur corps $k$

$v,w$ deux places archimédiennes de \deleted{$k$} $X$, d'où foncteurs fibres $\unsure{\zeta|\xi}^{\unsure{\mathbb{Z}|Z}}_v,\zeta^{\mathbb{Z}}_w$. Question : \inserted{a) Quand} sont-ils isomorphes ? b) \deleted{Est-il vrai que} Quand est-il vrai que pour tte sous-catégorie \inserted{mult.} de t. f. de $M(X)$, les restrictions de ces foncteurs \unsure{à} $M_\alpha(X)$ sont isomorphes ?

Nous pouvons supposer, quitte à remplacer $k$ par \deleted{\unsure{Soient $v',w'$ les}} \unsure{sa|la} clôture alg. ds $k(X)$, que $k$ est alg. clos ds $k(X)$ [utiliser le fait $X$ \uline{normal} !] donc $X$ est \unsure{géom.\textsuperscript{t}} irréd. sur $k$. Soit \deleted{$k_0$} $k_0$ la clôture alg. de $\unsure{\mathbb{Q}|Q}$ dans $k$.

\deleted{\unsure{le corps de coïncidence des places} (\unsure{= valuations}). \unsure{induites} sur $k$,}

\deleted{\unsure{envisager}}

\inserted{\uline{Contient la clôture alg. \deleted{$E$} de $\mathbb{Q}$ dans $k$} (p. ex. $k_0$ alg. fermé dans $k$).}

\deleted{1) $k_0$ \uline{est algébriquement fermé dans $k$}. Alors}

1\textsuperscript{o}) $v,\unsure{w'|w}$ coïncident sur $k_0$.

\[
\text{a) }\zeta^{\mathbb{Z}}_v\mid M_\alpha(X)\simeq\zeta^{\mathbb{Z}}_w\mid M_\alpha(X)\qquad\text{pour }\forall\;\alpha.
\]

\diagram{\begin{tikzcd}[ampersand replacement=\&,column sep=small,row sep=small]
X_0\arrow[rr,dash]\arrow[d,dash] \& \& X\arrow[d,dash]\\
k_0\arrow[r,dash] \& \deleted{A_0}\arrow[r,dash] \& k
\end{tikzcd}}

\diagram{\begin{tikzcd}[ampersand replacement=\&,column sep=small,row sep=small,execute at end picture={
\draw[thin] ($(midbase.west)+(-.15,0)$) -- ($(highx.east)+(.1,.12)$);
\draw[thin] (midbase.west) -- ($(highx.east)+(.24,.13)$);
\draw[thin] (midbase.center) -- ($(highx.east)+(.38,.14)$);
\draw[thin] (midbase.east) -- ($(highx.east)+(.52,.15)$);
\draw[thin] ($(midbase.east)+(.15,0)$) -- ($(highx.east)+(.66,.16)$);
\draw[thin] ($(midbase.east)+(.30,0)$) -- ($(highx.east)+(.8,.14)$);
\draw[thin] (lowq.west) .. controls +(0,-.4) and +(0,-.4) .. (lowe.east) .. controls +(.8,.9) and +(.7,.55) .. (lowe.north);
}]
\& \& \deleted{\unsure{X_0}}\arrow[rr,dash] \& \& X\arrow[dd,dash]\\
\& \& |[alias=highx]|\deleted{X_0}\arrow[d,dash]\arrow[urr,dash] \& \& \\
E\arrow[r,dash]\arrow[drr,dash] \& \unsure{k_0}\arrow[r,dash]\arrow[d,dash] \& |[alias=midbase]|\deleted{\unsure{A_0}}\arrow[rr,dash]\arrow[d,dash] \& \& k\\
\& |[alias=lowq]|\deleted{\unsure{Q|\mathbb{Q}}}\arrow[rr,dash] \& A_0 \& |[alias=lowe]|\deleted{\unsure{E|\mathbb{Z}}} \&
\end{tikzcd}}

En effet, $X$ et $M_\alpha(X)$ \unsure{provient|proviennent} de $X_0$, \inserted{$M_\alpha(X_0)$} de type fini sur $A_0$, \inserted{$k_0$} \unsure{ss alg.|alg.} de type fini de $k$, /\deleted{\unsure{et}} \deleted{\unsure{$X_0\to\operatorname{Sp} A_0$ fid plat et univ.|$X_0\to\operatorname{Spec} A_0$ fid plat et univ.}, \unsure{$\operatorname{Sp} A_0\to k$}} On est ramené au cas où $k=\deleted{\unsure{E}}k_0$, où on obtient le résultat voulu grâce au groupe fondamental transcendant. \deleted{\unsure{Mais}}

b) Si \deleted{\unsure{$k\subset k_0$}} \inserted{$v$ et $w$ coïncident sur $k$,} cet argument montre que $\zeta_v^{\mathbb{Z}}\simeq\zeta_w^{\mathbb{Z}}$. \deleted{Si \unsure{$k\ne k_0$},} \inserted{Sinon,} \deleted{je pense que $\zeta_v^{\mathbb{Z}}\ne\zeta_w^{\mathbb{Z}}$, \unsure{du moins}} \inserted{mais \unsure{si on suppose que} $k$ est} \deleted{\unsure{de type}}

\margin{/\;on peut supposer $X_0\to\operatorname{Spec} A_0$ fid plat \unsure{univers\textsuperscript{t}} connexe, et $\operatorname{Spec}(A_0)\mathbin{-} k_0$ lisse connexe, donc $X_0\mathbin{-}\operatorname{Spec} k_0$ lisse connexe.}
